\editionPage[\pdfbookmark[0]{Défis d’approfondissement}{bm-defis-debut-parent}]{defis-debut}{Vingt-et-un défis}{Un défi par chapitre · Parcours facultatif}{Chercher avant de consulter les deux indications.}
\begin{blocEdition}{defis-debut}
Ces problèmes demandent davantage d’initiative. Ils utilisent les outils du livre ; les prolongements sont définis dans l’énoncé. Les défis sont facultatifs et approfondissent les notions du programme.\par\bigskip
\begin{savoirfaire}[Une recherche qui avance]Faire un dessin, tester un petit cas, organiser des essais dans un tableau, garder ses traces. Ouvrir la première indication seulement après une vraie recherche ; revenir au problème avant de lire la seconde.\end{savoirfaire}
\begin{savoirfaire}[Après la solution]Repérer l’idée décisive et la refaire sur une feuille vierge, puis raconter la solution à quelqu’un. Expliquer pourquoi cela fonctionne compte autant que trouver.\end{savoirfaire}
\end{blocEdition}
\editionPage{defi-ex-1}{D01 · Défi}{Les nombres miroirs}{Approfondissement · 20 à 45 minutes indicatives}
\begin{blocEdition}{defi-ex-1}
Un nombre miroir se lit de la même façon de gauche à droite et de droite à gauche, comme $43\,234$. Combien y a-t-il de nombres miroirs à cinq chiffres compris entre $40\,000$ et $50\,000$ ? Donner le plus petit et le plus grand.\par\vspace{5mm}{\otLight\fontsize{8}{10}\selectfont\color{otMuted}Espace de recherche}\par\par\vspace{7mm}{\color{otRule!65}\hrule height .25pt}\par\vspace{7mm}{\color{otRule!65}\hrule height .25pt}\par\vspace{7mm}{\color{otRule!65}\hrule height .25pt}\par\vspace{7mm}{\color{otRule!65}\hrule height .25pt}\par\vspace{7mm}{\color{otRule!65}\hrule height .25pt}\par\vspace{7mm}{\color{otRule!65}\hrule height .25pt}
\end{blocEdition}
\editionPage{defi-ex-2}{D02 · Défi}{Somme et écart}{Approfondissement · 20 à 45 minutes indicatives}
\begin{blocEdition}{defi-ex-2}
Deux nombres ont pour somme $1\,000$ et leur écart vaut 248. Trouver ces deux nombres. Représenter la situation par un schéma en barres.\par\vspace{5mm}{\otLight\fontsize{8}{10}\selectfont\color{otMuted}Espace de recherche}\par\par\vspace{7mm}{\color{otRule!65}\hrule height .25pt}\par\vspace{7mm}{\color{otRule!65}\hrule height .25pt}\par\vspace{7mm}{\color{otRule!65}\hrule height .25pt}\par\vspace{7mm}{\color{otRule!65}\hrule height .25pt}\par\vspace{7mm}{\color{otRule!65}\hrule height .25pt}\par\vspace{7mm}{\color{otRule!65}\hrule height .25pt}
\end{blocEdition}
\editionPage{defi-ex-3}{D03 · Défi}{Les élèves en rangs}{Approfondissement · 20 à 45 minutes indicatives}
\begin{blocEdition}{defi-ex-3}
Un nombre d’élèves compris entre 100 et 200 se range en rangs. Par rangs de 4, il en reste 3 ; par rangs de 5, il en reste 4 ; par rangs de 6, il en reste 5. Combien d’élèves peut-il y avoir ? Trouver toutes les possibilités.\par\vspace{5mm}{\otLight\fontsize{8}{10}\selectfont\color{otMuted}Espace de recherche}\par\par\vspace{7mm}{\color{otRule!65}\hrule height .25pt}\par\vspace{7mm}{\color{otRule!65}\hrule height .25pt}\par\vspace{7mm}{\color{otRule!65}\hrule height .25pt}\par\vspace{7mm}{\color{otRule!65}\hrule height .25pt}\par\vspace{7mm}{\color{otRule!65}\hrule height .25pt}\par\vspace{7mm}{\color{otRule!65}\hrule height .25pt}
\end{blocEdition}
\editionPage{defi-ex-4}{D04 · Défi}{Les divisions qui s’arrêtent}{Approfondissement · 20 à 45 minutes indicatives}
\begin{blocEdition}{defi-ex-4}
On divise 1 successivement par 2, 3, 4, 5, 6, 7, 8 et 9, en poursuivant après la virgule. Quelles divisions s’arrêtent, avec un reste nul ? Donner leur quotient exact.\par\vspace{5mm}{\otLight\fontsize{8}{10}\selectfont\color{otMuted}Espace de recherche}\par\par\vspace{7mm}{\color{otRule!65}\hrule height .25pt}\par\vspace{7mm}{\color{otRule!65}\hrule height .25pt}\par\vspace{7mm}{\color{otRule!65}\hrule height .25pt}\par\vspace{7mm}{\color{otRule!65}\hrule height .25pt}\par\vspace{7mm}{\color{otRule!65}\hrule height .25pt}\par\vspace{7mm}{\color{otRule!65}\hrule height .25pt}
\end{blocEdition}
\editionPage{defi-ex-5}{D05 · Défi}{La suite des lapins}{Approfondissement · 20 à 45 minutes indicatives}
\begin{blocEdition}{defi-ex-5}
Dans la suite $1 ;\ 1 ;\ 2 ;\ 3 ;\ 5 ;\ 8 ;\ \ldots$ chaque terme est la somme des deux précédents. Donner le douzième terme. Combien de termes sont inférieurs à $1\,000$ ?\par\vspace{5mm}{\otLight\fontsize{8}{10}\selectfont\color{otMuted}Espace de recherche}\par\par\vspace{7mm}{\color{otRule!65}\hrule height .25pt}\par\vspace{7mm}{\color{otRule!65}\hrule height .25pt}\par\vspace{7mm}{\color{otRule!65}\hrule height .25pt}\par\vspace{7mm}{\color{otRule!65}\hrule height .25pt}\par\vspace{7mm}{\color{otRule!65}\hrule height .25pt}\par\vspace{7mm}{\color{otRule!65}\hrule height .25pt}
\end{blocEdition}
\editionPage{defi-ex-6}{D06 · Défi}{La fourmi pressée}{Approfondissement · 20 à 45 minutes indicatives}
\begin{blocEdition}{defi-ex-6}
Une fourmi part de 0 sur une règle d’un mètre. À chaque saut, elle parcourt la moitié de la distance qui lui reste jusqu’au bout. Où est-elle après 1, 2, 3, 4 et 5 sauts ? Quelle distance lui reste-t-il après 5 sauts ?\par\vspace{5mm}{\otLight\fontsize{8}{10}\selectfont\color{otMuted}Espace de recherche}\par\par\vspace{7mm}{\color{otRule!65}\hrule height .25pt}\par\vspace{7mm}{\color{otRule!65}\hrule height .25pt}\par\vspace{7mm}{\color{otRule!65}\hrule height .25pt}\par\vspace{7mm}{\color{otRule!65}\hrule height .25pt}\par\vspace{7mm}{\color{otRule!65}\hrule height .25pt}\par\vspace{7mm}{\color{otRule!65}\hrule height .25pt}
\end{blocEdition}
\editionPage{defi-ex-7}{D07 · Défi}{Les douzièmes du milieu}{Approfondissement · 20 à 45 minutes indicatives}
\begin{blocEdition}{defi-ex-7}
Trouver toutes les fractions de dénominateur 12 strictement comprises entre $\frac13$ et $\frac34$. Laquelle est égale à un demi ?\par\vspace{5mm}{\otLight\fontsize{8}{10}\selectfont\color{otMuted}Espace de recherche}\par\par\vspace{7mm}{\color{otRule!65}\hrule height .25pt}\par\vspace{7mm}{\color{otRule!65}\hrule height .25pt}\par\vspace{7mm}{\color{otRule!65}\hrule height .25pt}\par\vspace{7mm}{\color{otRule!65}\hrule height .25pt}\par\vspace{7mm}{\color{otRule!65}\hrule height .25pt}\par\vspace{7mm}{\color{otRule!65}\hrule height .25pt}
\end{blocEdition}
\editionPage{defi-ex-8}{D08 · Défi}{L’héritage des chameaux}{Approfondissement · 20 à 45 minutes indicatives}
\begin{blocEdition}{defi-ex-8}
Un marchand laisse 17 chameaux à ses trois enfants : l’aîné doit recevoir la moitié, le second le tiers et le dernier le neuvième. Personne ne veut couper de chameau ! Un voyageur prête son propre chameau, partage les 18 chameaux, puis repart avec le sien. Comment ? Pourquoi cela fonctionne-t-il ?\par\vspace{5mm}{\otLight\fontsize{8}{10}\selectfont\color{otMuted}Espace de recherche}\par\par\vspace{7mm}{\color{otRule!65}\hrule height .25pt}\par\vspace{7mm}{\color{otRule!65}\hrule height .25pt}\par\vspace{7mm}{\color{otRule!65}\hrule height .25pt}\par\vspace{7mm}{\color{otRule!65}\hrule height .25pt}\par\vspace{7mm}{\color{otRule!65}\hrule height .25pt}\par\vspace{7mm}{\color{otRule!65}\hrule height .25pt}
\end{blocEdition}
\editionPage{defi-ex-9}{D09 · Défi}{Le décimal à code}{Approfondissement · 20 à 45 minutes indicatives}
\begin{blocEdition}{defi-ex-9}
Je suis un nombre compris entre 3 et 4. Mon chiffre des centièmes est 0, mon chiffre des millièmes est le double de mon chiffre des dixièmes, et la somme de mes chiffres après la virgule vaut 9. Qui suis-je ?\par\vspace{5mm}{\otLight\fontsize{8}{10}\selectfont\color{otMuted}Espace de recherche}\par\par\vspace{7mm}{\color{otRule!65}\hrule height .25pt}\par\vspace{7mm}{\color{otRule!65}\hrule height .25pt}\par\vspace{7mm}{\color{otRule!65}\hrule height .25pt}\par\vspace{7mm}{\color{otRule!65}\hrule height .25pt}\par\vspace{7mm}{\color{otRule!65}\hrule height .25pt}\par\vspace{7mm}{\color{otRule!65}\hrule height .25pt}
\end{blocEdition}
\editionPage{defi-ex-10}{D10 · Défi}{La tirelire à deux pièces}{Approfondissement · 20 à 45 minutes indicatives}
\begin{blocEdition}{defi-ex-10}
Une tirelire contient exactement 10 pièces, uniquement des pièces de 50 centimes et de 20 centimes, pour un total de $3{,}50$ euros. Combien y a-t-il de pièces de chaque sorte ?\par\vspace{5mm}{\otLight\fontsize{8}{10}\selectfont\color{otMuted}Espace de recherche}\par\par\vspace{7mm}{\color{otRule!65}\hrule height .25pt}\par\vspace{7mm}{\color{otRule!65}\hrule height .25pt}\par\vspace{7mm}{\color{otRule!65}\hrule height .25pt}\par\vspace{7mm}{\color{otRule!65}\hrule height .25pt}\par\vspace{7mm}{\color{otRule!65}\hrule height .25pt}\par\vspace{7mm}{\color{otRule!65}\hrule height .25pt}
\end{blocEdition}
\editionPage{defi-ex-11}{D11 · Défi}{La balance aux trois masses}{Approfondissement · 20 à 45 minutes indicatives}
\begin{blocEdition}{defi-ex-11}
On dispose d’une balance à deux plateaux et de trois masses : 1 g, 3 g et 9 g. On peut poser des masses sur les deux plateaux. Montrer qu’on peut peser tout objet de 1 g à 13 g. Donner la méthode pour 2 g, 5 g et 11 g.\par\vspace{5mm}{\otLight\fontsize{8}{10}\selectfont\color{otMuted}Espace de recherche}\par\par\vspace{7mm}{\color{otRule!65}\hrule height .25pt}\par\vspace{7mm}{\color{otRule!65}\hrule height .25pt}\par\vspace{7mm}{\color{otRule!65}\hrule height .25pt}\par\vspace{7mm}{\color{otRule!65}\hrule height .25pt}\par\vspace{7mm}{\color{otRule!65}\hrule height .25pt}\par\vspace{7mm}{\color{otRule!65}\hrule height .25pt}
\end{blocEdition}
\editionPage{defi-ex-12}{D12 · Défi}{Le plus grand enclos}{Approfondissement · 20 à 45 minutes indicatives}
\begin{blocEdition}{defi-ex-12}
Avec 24 m de grillage, on veut entourer un enclos rectangulaire dont les côtés mesurent un nombre entier de mètres. Lister tous les rectangles possibles et trouver celui qui a la plus grande aire.\par\vspace{5mm}{\otLight\fontsize{8}{10}\selectfont\color{otMuted}Espace de recherche}\par\par\vspace{7mm}{\color{otRule!65}\hrule height .25pt}\par\vspace{7mm}{\color{otRule!65}\hrule height .25pt}\par\vspace{7mm}{\color{otRule!65}\hrule height .25pt}\par\vspace{7mm}{\color{otRule!65}\hrule height .25pt}\par\vspace{7mm}{\color{otRule!65}\hrule height .25pt}\par\vspace{7mm}{\color{otRule!65}\hrule height .25pt}
\end{blocEdition}
\editionPage{defi-ex-13}{D13 · Défi}{L’horloge des angles}{Approfondissement · 20 à 45 minutes indicatives}
\begin{blocEdition}{defi-ex-13}
À 3 h pile, les aiguilles d’une horloge forment un angle droit. Sachant que la grande aiguille indique 12 à chaque heure pile, quel angle forment les aiguilles à 2 h ? à 4 h ? À quelles heures piles l’angle mesure-t-il 150 degrés ?\par\vspace{5mm}{\otLight\fontsize{8}{10}\selectfont\color{otMuted}Espace de recherche}\par\par\vspace{7mm}{\color{otRule!65}\hrule height .25pt}\par\vspace{7mm}{\color{otRule!65}\hrule height .25pt}\par\vspace{7mm}{\color{otRule!65}\hrule height .25pt}\par\vspace{7mm}{\color{otRule!65}\hrule height .25pt}\par\vspace{7mm}{\color{otRule!65}\hrule height .25pt}\par\vspace{7mm}{\color{otRule!65}\hrule height .25pt}
\end{blocEdition}
\editionPage{defi-ex-14}{D14 · Défi}{Les deux sabliers}{Approfondissement · 20 à 45 minutes indicatives}
\begin{blocEdition}{defi-ex-14}
On dispose d’un sablier de 7 minutes et d’un sablier de 4 minutes. Comment mesurer exactement 9 minutes ?\par\vspace{5mm}{\otLight\fontsize{8}{10}\selectfont\color{otMuted}Espace de recherche}\par\par\vspace{7mm}{\color{otRule!65}\hrule height .25pt}\par\vspace{7mm}{\color{otRule!65}\hrule height .25pt}\par\vspace{7mm}{\color{otRule!65}\hrule height .25pt}\par\vspace{7mm}{\color{otRule!65}\hrule height .25pt}\par\vspace{7mm}{\color{otRule!65}\hrule height .25pt}\par\vspace{7mm}{\color{otRule!65}\hrule height .25pt}
\end{blocEdition}
\editionPage{defi-ex-15}{D15 · Défi}{Le trésor à deux distances}{Approfondissement · 20 à 45 minutes indicatives}
\begin{blocEdition}{defi-ex-15}
Le puits P et le chêne C sont à 5 cm l’un de l’autre sur le plan. Le trésor est à 4 cm du puits et à 3 cm du chêne. Construire les positions possibles du trésor. Combien y en a-t-il ?\par\vspace{5mm}{\otLight\fontsize{8}{10}\selectfont\color{otMuted}Espace de recherche}\par\par\vspace{7mm}{\color{otRule!65}\hrule height .25pt}\par\vspace{7mm}{\color{otRule!65}\hrule height .25pt}\par\vspace{7mm}{\color{otRule!65}\hrule height .25pt}\par\vspace{7mm}{\color{otRule!65}\hrule height .25pt}\par\vspace{7mm}{\color{otRule!65}\hrule height .25pt}\par\vspace{7mm}{\color{otRule!65}\hrule height .25pt}
\end{blocEdition}
\editionPage{defi-ex-16}{D16 · Défi}{Les diagonales des polygones}{Approfondissement · 20 à 45 minutes indicatives}
\begin{blocEdition}{defi-ex-16}
Une diagonale relie deux sommets non voisins d’un polygone. Combien de diagonales possède un quadrilatère ? un pentagone ? un hexagone ?\par\vspace{5mm}{\otLight\fontsize{8}{10}\selectfont\color{otMuted}Espace de recherche}\par\par\vspace{7mm}{\color{otRule!65}\hrule height .25pt}\par\vspace{7mm}{\color{otRule!65}\hrule height .25pt}\par\vspace{7mm}{\color{otRule!65}\hrule height .25pt}\par\vspace{7mm}{\color{otRule!65}\hrule height .25pt}\par\vspace{7mm}{\color{otRule!65}\hrule height .25pt}\par\vspace{7mm}{\color{otRule!65}\hrule height .25pt}
\end{blocEdition}
\editionPage{defi-ex-17}{D17 · Défi}{Les mots au fil de l’eau}{Approfondissement · 20 à 45 minutes indicatives}
\begin{blocEdition}{defi-ex-17}
Les lettres B, C, D, E, H, I, K, O et X ont un axe de symétrie horizontal. Expliquer pourquoi le mot CODE se lit encore quand on le reflète dans un lac, axe horizontal. Trouver un autre mot qui résiste à ce reflet.\par\vspace{5mm}{\otLight\fontsize{8}{10}\selectfont\color{otMuted}Espace de recherche}\par\par\vspace{7mm}{\color{otRule!65}\hrule height .25pt}\par\vspace{7mm}{\color{otRule!65}\hrule height .25pt}\par\vspace{7mm}{\color{otRule!65}\hrule height .25pt}\par\vspace{7mm}{\color{otRule!65}\hrule height .25pt}\par\vspace{7mm}{\color{otRule!65}\hrule height .25pt}\par\vspace{7mm}{\color{otRule!65}\hrule height .25pt}
\end{blocEdition}
\editionPage{defi-ex-18}{D18 · Défi}{Le grand cube peint}{Approfondissement · 20 à 45 minutes indicatives}
\begin{blocEdition}{defi-ex-18}
Un grand cube est formé de 64 petits cubes, 4 sur 4 sur 4. On peint toute sa surface. Combien de petits cubes ont 3 faces peintes ? 2 faces ? 1 face ? aucune ?\par\vspace{5mm}{\otLight\fontsize{8}{10}\selectfont\color{otMuted}Espace de recherche}\par\par\vspace{7mm}{\color{otRule!65}\hrule height .25pt}\par\vspace{7mm}{\color{otRule!65}\hrule height .25pt}\par\vspace{7mm}{\color{otRule!65}\hrule height .25pt}\par\vspace{7mm}{\color{otRule!65}\hrule height .25pt}\par\vspace{7mm}{\color{otRule!65}\hrule height .25pt}\par\vspace{7mm}{\color{otRule!65}\hrule height .25pt}
\end{blocEdition}
\editionPage{defi-ex-19}{D19 · Défi}{Deux robinets ensemble}{Approfondissement · 20 à 45 minutes indicatives}
\begin{blocEdition}{defi-ex-19}
Le robinet A remplit 30 L en 5 minutes ; le robinet B remplit 24 L en 6 minutes ; chacun coule régulièrement. Ouverts ensemble, en combien de minutes remplissent-ils une citerne de 270 L ?\par\vspace{5mm}{\otLight\fontsize{8}{10}\selectfont\color{otMuted}Espace de recherche}\par\par\vspace{7mm}{\color{otRule!65}\hrule height .25pt}\par\vspace{7mm}{\color{otRule!65}\hrule height .25pt}\par\vspace{7mm}{\color{otRule!65}\hrule height .25pt}\par\vspace{7mm}{\color{otRule!65}\hrule height .25pt}\par\vspace{7mm}{\color{otRule!65}\hrule height .25pt}\par\vspace{7mm}{\color{otRule!65}\hrule height .25pt}
\end{blocEdition}
\editionPage{defi-ex-20}{D20 · Défi}{La récréation en secteurs}{Approfondissement · 20 à 45 minutes indicatives}
\begin{blocEdition}{defi-ex-20}
Pendant la récréation, 120 élèves jouent : la moitié au ballon, le tiers à la corde, et les autres aux billes. Combien d’élèves jouent aux billes ? Quelle fraction du disque occupe le secteur des billes dans un diagramme circulaire ?\par\vspace{5mm}{\otLight\fontsize{8}{10}\selectfont\color{otMuted}Espace de recherche}\par\par\vspace{7mm}{\color{otRule!65}\hrule height .25pt}\par\vspace{7mm}{\color{otRule!65}\hrule height .25pt}\par\vspace{7mm}{\color{otRule!65}\hrule height .25pt}\par\vspace{7mm}{\color{otRule!65}\hrule height .25pt}\par\vspace{7mm}{\color{otRule!65}\hrule height .25pt}\par\vspace{7mm}{\color{otRule!65}\hrule height .25pt}
\end{blocEdition}
\editionPage{defi-ex-21}{D21 · Défi}{La somme la plus probable}{Approfondissement · 20 à 45 minutes indicatives}
\begin{blocEdition}{defi-ex-21}
On lance deux dés équilibrés et on additionne les deux faces. Quelle somme a le plus de chances de sortir ? Donner ses chances sous la forme « a chances sur b ». La somme 2 et la somme 12 ont-elles les mêmes chances ?\par\vspace{5mm}{\otLight\fontsize{8}{10}\selectfont\color{otMuted}Espace de recherche}\par\par\vspace{7mm}{\color{otRule!65}\hrule height .25pt}\par\vspace{7mm}{\color{otRule!65}\hrule height .25pt}\par\vspace{7mm}{\color{otRule!65}\hrule height .25pt}\par\vspace{7mm}{\color{otRule!65}\hrule height .25pt}\par\vspace{7mm}{\color{otRule!65}\hrule height .25pt}\par\vspace{7mm}{\color{otRule!65}\hrule height .25pt}
\end{blocEdition}
